\cleardoublepage

\section{方法论与设计思路}

\subsection{系统架构概述}
本文提出的 VLN-R1 系统采用了一种多阶段微调与层次化规划相结合的架构，旨在解决真实机器人在长程导航任务中的逻辑推理与动作一致性问题。系统整体由视觉感知引擎、拓扑建图模块、跨模态推理模块以及底层控制器四个核心部分组成（如图\ref{fig:framework}所示，注：图引用自开题报告）。

系统的任务处理流程如下：
1. **多模态感知**：利用 DINOv2 视觉编码器提取当前视角的语义特征，并结合坐标系对齐后的四目相机（Seeker）图像进行特征融合。
2. **拓扑图构建**：基于路点预测器在线构建环境拓扑图，将复杂的三维空间简化为一系列可通行的节点与边。
3. **思维链推理（CoT）**：跨模态规划器在决策前产生一段思维链推理（Think），分析当前观测、历史指令与拓扑结构，生成下一步的宏观导航目标。
4. **动作解码与执行**：将宏观目标转化为底层连续控制指令（$v, \omega$），通过闭环控制实现平稳移动。

\subsection{感知迁移与 4-View 特征适配}
针对仿真环境与实车硬件之间的视场角（FoV）差异，本系统设计了特定的特征适配层。
\subsubsection{DINOv2 视觉编码器集成}
相比传统的 CLIP 编码器，DINOv2 通过自监督学习获取了更强的几何深度一致性特征。我们在模型的前端引入了 DINOv2-ViT-B/14 骨干网络，通过对不同视角的图像进行特征提取，并利用线性投影层将其映射至具身智能体的语义嵌入空间。

\subsubsection{Seeker 坐标系对齐算法}
仿真器通常提供 $360^{\circ}$ 的全景（12-View）视野，而实车 Seeker 相机仅提供正面约 $200^{\circ}$ 的 4-View 重叠视野。本系统通过以下步骤实现视场对齐：
1. **相机内参校准**：获取实车 4 台相机的焦距与主点坐标。
2. **重投影变换**：利用旋转矩阵 $\mathbf{R}$ 将 12 目仿真的观测采样至与 Seeker 相机视场一致的视角。
3. **多视角注意力融合**：在输入模型前，利用自注意力机制对 4 个视角的视觉信息进行加权融合，使模型能够自适应地关注关键路标。

\subsection{VLN-Ego: CoT 增强数据集构建}
为了赋予模型推理能力，我们设计了自动化 CoT 数据引擎。该引擎基于 RxR 数据集的原始轨迹，利用 Qwen3-VL-32B 这种具备深度思考能力的多模态大模型，为每一个决策步骤生成对应的“思维过程”。

数据集定义为四元组序列 $\mathcal{D} = \{ (O_t, I, C_t, A_t) \}_{t=1}^T$，其中：
- $O_t$ 为当前的多维视角观测。
- $I$ 为原始自然语言导航指令。
- $C_t$ 是思维链推理内容（如：“我正处于走廊入口，左侧是一个开放式厨房，我需要按照指令继续直行以寻找卧室”）。
- $A_t$ 是对应的控制动作。

\subsection{多阶段强化微调策略}
本系统放弃了直接从随机状态训练的方式，采用了“SFT 冷启动 + GRPO 提升”的训练模式。

\subsubsection{TBA-SFT 冷启动阶段}
在第一阶段，模型通过监督微调（Supervised Fine-Tuning）学习基础的格式规范与导航逻辑。我们通过 VLN-Ego 数据集，强迫模型输出符合 \texttt{<Think>...</Think><Action>...</Action>} 格式的内容。这一阶段初步打通了感知与动作之间的语义映射。

\subsubsection{基于 GRPO 的策略优化}
在具备基础能力后，我们引入组相对策略优化算法进行端到端微调。
\textbf{奖励函数设计}：
我们设计了三维复合奖励以约束模型的行为：
\begin{itemize}
    \item \textbf{格式奖励（$R_{format}$）}：检查输出标签是否完整，推理内容是否冗余。
    \item \textbf{语义理解奖励（$R_{sim}$）}：利用预训练的 CLIP 计算模型生成的“思维内容”与当前视觉场景的余弦相似度。
    \item \textbf{导航成功奖励（$R_{nav}$）}：基于 SPL 指标和终点误差进行即时反馈。
\end{itemize}

通过 GRPO 算法，我们在每一轮采样 $G$ 条轨迹，根据组内的相对优势更新梯度，避免了复杂价值函数的误差累积，使模型能够在长时程的任务中保持规划的一致性。
